% Official source: https://www.imo-official.org/assets/documents/problems/2013/2013_eng.pdf
\competitionyear{2013}

\begin{problem}{2013}{1}
Given positive integers $k, n$, prove that some positive integers $m_1, \dots, m_k$, not necessarily distinct, satisfy
\[
    1+\frac{2^k-1}{n}=\prod_{i=1}^k\left(1+\frac1{m_i}\right).
\]
\end{problem}

\begin{problem}{2013}{2}
A configuration of $4027$ points in the plane is called \emph{Colombian} if $2013$ of its points are red and $2014$ are blue, with no three collinear. A collection of lines is \emph{good} for such a configuration if
\begin{problemchoices}
    \item none of its lines passes through a point of the configuration;
    \item every region into which the lines divide the plane contains points of at most one colour.
\end{problemchoices}
Find the smallest $k$ such that every Colombian configuration admits a good collection of $k$ lines.
\end{problem}

\begin{problem}{2013}{3}
Let the excircle of triangle $ABC$ opposite $A$ touch $BC$ at $A_1$. Let $B_1$ and $C_1$ be the analogous tangency points on the sides $CA$ and $AB$ for the excircles opposite $B$ and $C$. Suppose that the circumcentre of $A_1B_1C_1$ lies on the circumcircle of $ABC$. Prove that $ABC$ is a right triangle.

Here the excircle opposite $A$ touches the segment $BC$ and the rays $AB$, $AC$ beyond $B$, $C$, respectively; the other excircles are defined analogously.
\end{problem}

\begin{problem}{2013}{4}
Let $ABC$ be an acute triangle with orthocentre $H$, and let $W$ lie strictly between $B$ and $C$ on the side $BC$. Let $M$ and $N$ be the feet of the altitudes from $B$ and $C$, respectively. Let $\omega_1$ be the circumcircle of triangle $BWN$, and let $X\in\omega_1$ be such that $WX$ is a diameter. Likewise, let $\omega_2$ be the circumcircle of $CWM$, and let $Y\in\omega_2$ be such that $WY$ is a diameter. Prove that $X, Y, H$ are collinear.
\end{problem}

\begin{problem}{2013}{5}
Suppose that $f\colon\Q_{>0}\to\R$ satisfies
\begin{problemclauses}
    \item $f(x)f(y)\ge f(xy)$ for all $x, y\in\Q_{>0}$;
    \item $f(x+y)\ge f(x)+f(y)$ for all $x, y\in\Q_{>0}$;
    \item $f(a)=a$ for some rational number $a>1$.
\end{problemclauses}
Prove that $f(x)=x$ for every $x\in\Q_{>0}$.
\end{problem}

\begin{problem}{2013}{6}
Mark $n+1$ equally spaced points on a circle, where $n\ge3$. Label them with $0, 1, \dots, n$, each used once, and consider two labellings equivalent if one is obtained from the other by rotating the circle. A labelling is \emph{beautiful} if, whenever $a<b<c<d$ and $a+d=b+c$, the chord joining labels $a$ and $d$ does not intersect the chord joining labels $b$ and $c$.

Let $M$ count the beautiful labellings up to rotation, and let $N$ count ordered pairs $(x, y)$ of positive integers with $x+y\le n$ and $\gcd(x, y)=1$. Prove that $M=N+1$.
\end{problem}
